\section{対角混合族の完全な同型条件と反対代数}
\label{sec:pair-diagonal-classification}
$D=x_1^3+x_2^3+x_3^3$、$T$を六つの混合語の和とし、
$s=(u,v)^t$, $q_i\in k[u,v]_2$, $f\in k[u,v]_3$に対して
\begin{gather*}
 B(h,q)=\frac{k[u,v]\langle x_1,x_2,x_3\rangle}
 {([x_j,x_k]-h(x_i^2+q_i))_{(i,j,k)\text{ cyclic}}},\\
 W_{h,q,f}=D/3-h^2T/9+\sum_iq_ix_i+f,
 \quad A(h,q,f)=B(h,q)/(W_{h,q,f})
\end{gather*}
と置く。以下$h\ne0$とする。境界の滑らかさは同型判定には不要である。

\begin{lemma}
$Z(A)_1=ku\oplus kv$であり、$\Sym^2(A_1)\to A_2$は同型である。
\end{lemma}
\begin{proof}
二次関係は$\wedge^2A_1$から$\Sym^2A_1$へのgraphで、三次商は次数二を変えない。
従って対称二次式は$A_2$の基底である。
一次元$\sum a_ix_i+bu+cv$の交換子を取ると、独立な$x_i^2$の係数と$h\ne0$から
$a_i=0$を得る。
\end{proof}
\begin{theorem}[対角族の同型の必要十分条件]
\label{thm:pair-diagonal-isomorphism}
$h,h'\ne0$について$A(h',q',f')\simeq_{\gr} A(h,q,f)$であることは、
$p\in S_3$, $N\in\GL_2$, $\lambda_i\ne0$, $\kappa\ne0$が存在して
\begin{gather*}
 \lambda_1^3=\lambda_2^3=\lambda_3^3=\kappa,\qquad
 h'=\operatorname{sgn}(p)h\lambda_1\lambda_2\lambda_3/\kappa,\\
 q_i'(Ns)=\lambda_i^2q_{p(i)}(s),\qquad f'(Ns)=\kappa f(s)
\end{gather*}
を満たすことと同値である。全ての同型は
$s'\mapsto Ns$, $x_i'\mapsto\lambda_i x_{p(i)}$であり、特に$h^6$は不変量である。
\end{theorem}
\begin{proof}
二次包絡の復元により、同型は周囲へ持ち上がる。
中心一次部分から$s'\mapsto Ns$, $x'\mapsto Mx+\ell(s)$であり、
関係式を対称二次式として比較すると
\[
 h\sum_a\operatorname{cof}(M)_{ia}(x_a^2+q_a(s))
 =h'\bigl((Mx)_i+\ell_i(s)\bigr)^2+h'q_i'(Ns).
\]
$x_as_b$の係数から$\ell_i=0$、$x_ax_b$ ($a\ne b$)の係数から
$M_{ia}M_{ib}=0$となる。可逆性により$M$はmonomialである。
$x_{p(i)}^2$の係数を比較して
$h'\lambda_i^2=\operatorname{sgn}(p)h\lambda_j\lambda_k$を得て、
$\lambda_i$を掛ければ三つの立方は等しい。
中央二次部分の比較が$q_i$の条件を与える。
$L=\prod\lambda_i$と置けば$L^3=\kappa^3$, $(h')^2L=\kappa h^2$なので、
$W'$の像は$\kappa(D/3-h^2T/9+\sum q_ix_i)+f'(Ns)$となる。
最高active三次部分の比較から比例係数は$\kappa$であり、$f$の条件を得る。
この最高次元は$h^6=-27$でも非零である。
実際、三変数代数の可換化は$k[x,y,z]/(x^2,y^2,z^2)$で、
$\Omega_h$の像は$-2h^2xyz/3\ne0$である。
逆に表示条件を満たす可逆線形写像は全ての関係を保つので同型を与える。
$h'$の式の六乗が最後の不変量を与える。
\end{proof}

\begin{proposition}[軌道次元と自己同型]
$q_1q_2q_3$の六零点が相異なる範囲で、周囲の対角族の同型類は6次元、
中心三次対の同型類は10次元である。全37次元周囲族と41次元対の中で、
それぞれの$\PGL_5$飽和像の余次元は7となる。
さらに六点集合の$\PGL_2$ stabilizerが自明で$f\ne0$なら
$\Aut_{\gr}A=k^\times$、$A\not\simeq_{\gr} A^{\rm op}$である。
\end{proposition}
\begin{proof}
連続な変換パラメータは$(c,N)\in\mathbb G_m\times\GL_2$の5次元であり、
全生成元の同時スカラーという1次元核を持つ。
六点のstabilizerは有限で、各候補の射影$N$を固定すると$q_i$の条件は
$N$のスカラーと$c$の比を有限個に制限する。
従って軌道は4次元で、周囲の10係数・対の14係数からそれぞれ4を引く。
一般の$\PGL_5$軌道は24次元なので飽和像は30次元・34次元であり、余次元7を得る。
stabilizerが自明なら$N=tI$, $p=1$で、$\lambda_i^2=t^2$と立方の条件から
全ての$\lambda_i$は同じ$\lambda=\pm t$である。
$f\ne0$が$\lambda^3=t^3$を強制し、自己同型は同時スカラーのみとなる。
反対代数は$A(-h,q,f)$であるが、同じ条件は$h'=h$を強制して矛盾する。
\end{proof}
\begin{corollary}[有理係数の非自己反対例と標識付き導来圏]
\label{cor:pair-not-opposite}
$q_1=uv$, $q_2=(u-v)(u-2v)$, $q_3=(u-4v)(u-9v)$とする。
六点集合のstabilizerは自明である。
$h=\pm1$について、両方のtailsが正則となる$f$の非空開集合が存在し、
その中の$f\ne0$では$A$と$A^{\rm op}$の間に構造対象をline objectのshiftへ送る
導来同値は存在しない。有理係数のそのような$f$はZariski稠密に存在する。
\end{corollary}
\begin{proof}
$\{0,\infty,1,2,4,9\}$の最初の三点の像を選ぶ120通りの有理射影行列を列挙し、
全六点を保つものを照合すると恒等変換だけである。
\texttt{verify\_central\_pair\_orbits.py}はこの厳密列挙を行う。
$h=\pm1$の境界の位数は補題\ref{CSBGlem:order-six}により6であり、
定理\ref{CSBGthm:global-open}を両周囲に適用する。
全ての$f$で相対PBW・中心正則商の自由性があり、最高次中心元は非零正則なので全仮定が揃う。
二つの非空開集合は交わり、$f=0$を除いても非空である。
Serre関手・偶数Veroneseによる復元定理と前命題から導来同値の不存在を得る。
有理点の稠密性は$\mathbb Q^4\subset\mathbb C^4$のZariski稠密性による。
構造対象を忘れた全ての導来同値についての主張ではない。
\end{proof}
